% !TEX program = xelatex
%
%Academic CV LaTeX Template (Chinese version)
% Based on: cv/my_cv.tex
%
% Compile with: xelatex cv_zh.tex
% NOTE: MUST be compiled with XeLaTeX (ctex Chinese support), NOT pdflatex.
%

\documentclass[a4paper,11pt]{article}

% Package imports
\usepackage[UTF8, fontset=none]{ctex} % fonts configured manually below (labmate's setup)
\setmainfont{texgyrepagella}[ % Palatino clone, loaded by filename from TeX Live (portable)
  Extension = .otf,
  UprightFont = *-regular,
  BoldFont = *-bold,
  ItalicFont = *-italic,
  BoldItalicFont = *-bolditalic,
]
% CJK font chain, identical to labmate's Chinese_version.tex:
% prefers Noto Serif CJK SC / Source Han / PingFang when installed;
% on Windows resolves to SimSun with AutoFakeBold (fake-bold Song, not SimHei) + AutoFakeSlant
\newcommand{\setchinesefont}{%
  \IfFontExistsTF{Noto Serif CJK SC}{\setCJKmainfont{Noto Serif CJK SC}}{%
  \IfFontExistsTF{Noto Sans CJK SC}{\setCJKmainfont{Noto Sans CJK SC}}{%
  \IfFontExistsTF{Source Han Serif SC}{\setCJKmainfont{Source Han Serif SC}}{%
  \IfFontExistsTF{Source Han Sans SC}{\setCJKmainfont{Source Han Sans SC}}{%
  \IfFontExistsTF{PingFang SC}{\setCJKmainfont{PingFang SC}}{%
  \IfFontExistsTF{Songti SC}{\setCJKmainfont{Songti SC}}{%
  \IfFontExistsTF{SimSun}{\setCJKmainfont[AutoFakeBold=2.5,AutoFakeSlant=0.2]{SimSun}}{%
  \IfFontExistsTF{Microsoft YaHei}{\setCJKmainfont[AutoFakeBold=2.5,AutoFakeSlant=0.2]{Microsoft YaHei}}{%
  }}}}}}}}
}
\setchinesefont
\usepackage{latexsym}
\usepackage{xcolor}
\usepackage{float}
\usepackage{ragged2e}
\usepackage[empty]{fullpage}
\usepackage{wrapfig}
\usepackage{lipsum}
\usepackage{tabularx}
\usepackage{titlesec}
\usepackage{geometry}
\usepackage{marvosym}
\usepackage{verbatim}
\usepackage{enumitem}
\usepackage{fancyhdr}
\usepackage{multicol}
\usepackage{graphicx}
\usepackage{fontawesome5}

% Color definitions
\definecolor{darkblue}{RGB}{0,0,139}

% Page layout
\setlength{\multicolsep}{0pt}
\pagestyle{fancy}
\fancyhf{} % clear all header and footer fields
\fancyfoot{}
\renewcommand{\headrulewidth}{0pt}
\renewcommand{\footrulewidth}{0pt}
\geometry{left=1.4cm, top=0.8cm, right=1.2cm, bottom=1cm}
\setlength{\footskip}{5pt} % Addressing fancyhdr warning

% Hyperlink setup (moved after fancyhdr to address warning)
\usepackage[hidelinks]{hyperref}
\hypersetup{
    colorlinks=true,
    linkcolor=darkblue,
    filecolor=darkblue,
    urlcolor=darkblue,
}

% Custom box settings
\usepackage[most]{tcolorbox}
\tcbset{
    frame code={},
    center title,
    left=0pt,
    right=0pt,
    top=0pt,
    bottom=0pt,
    colback=gray!20,
    colframe=white,
    width=\dimexpr\textwidth\relax,
    enlarge left by=-2mm,
    boxsep=4pt,
    arc=0pt,outer arc=0pt,
}

% URL style
\urlstyle{same}

% Text alignment
\raggedright
\setlength{\tabcolsep}{0in}

% Section formatting
\titleformat{\section}{
  \vspace{-4pt}\scshape\raggedright\large
}{}{0em}{}[\color{black}\titlerule \vspace{-7pt}]

% Custom commands
\newcommand{\resumeItem}[2]{
  \item{
    \textbf{#1}{\hspace{0.5mm}#2 \vspace{-0.5mm}}
  }
}

\newcommand{\resumePOR}[3]{
\vspace{0.5mm}\item
    \begin{tabular*}{0.97\textwidth}[t]{l@{\extracolsep{\fill}}r}
        \textbf{#1}\hspace{0.3mm}#2 & \textit{\small{#3}}
    \end{tabular*}
    \vspace{-2mm}
}

\newcommand{\resumeSubheading}[4]{
\vspace{0.5mm}\item
    \begin{tabular*}{0.98\textwidth}[t]{l@{\extracolsep{\fill}}r}
        \textbf{#1} & \textit{\footnotesize{#4}} \\
        \textit{\footnotesize{#3}} &  \footnotesize{#2}\\
    \end{tabular*}
    \vspace{-2.4mm}
}

\newcommand{\resumeProject}[4]{
\vspace{0.5mm}\item
    \begin{tabularx}{0.98\textwidth}[t]{@{}Xr@{}}
        \textbf{#1} & \textit{\footnotesize{#3}} \\
        \footnotesize{\textit{#2}} & \footnotesize{#4}
    \end{tabularx}
    \vspace{-2.4mm}
}

\newcommand{\resumeSubItem}[2]{\resumeItem{#1}{#2}\vspace{-4pt}}

\renewcommand{\labelitemi}{$\vcenter{\hbox{\tiny$\bullet$}}$}
\renewcommand{\labelitemii}{$\vcenter{\hbox{\tiny$\circ$}}$}

\newcommand{\resumeSubHeadingListStart}{\begin{itemize}[leftmargin=*,labelsep=1mm]}
\newcommand{\resumeHeadingSkillStart}{\begin{itemize}[leftmargin=*,itemsep=1.7mm, rightmargin=2ex]}
\newcommand{\resumeItemListStart}{\begin{itemize}[leftmargin=*,labelsep=1mm,itemsep=0mm]}

\newcommand{\resumeSubHeadingListEnd}{\end{itemize}\vspace{2mm}}
\newcommand{\resumeHeadingSkillEnd}{\end{itemize}\vspace{-2mm}}
\newcommand{\resumeItemListEnd}{\end{itemize}\vspace{-2mm}}
\newcommand{\cvsection}[1]{%
\vspace{2mm}
\begin{tcolorbox}
    \textbf{\large #1}
\end{tcolorbox}
    \vspace{-4mm}
}

\newcolumntype{L}{>{\raggedright\arraybackslash}X}%
\newcolumntype{R}{>{\raggedleft\arraybackslash}X}%
\newcolumntype{C}{>{\centering\arraybackslash}X}%

% Commands for icon sizing and positioning
\newcommand{\socialicon}[1]{\raisebox{-0.05em}{\resizebox{!}{1em}{#1}}}
\newcommand{\ieeeicon}[1]{\raisebox{-0.3em}{\resizebox{!}{1.3em}{#1}}}

% Font options (Latin text)
\newcommand{\headerfonti}{\fontfamily{phv}\selectfont} % Helvetica-like (similar to Arial/Calibri)
\newcommand{\headerfontii}{\fontfamily{ptm}\selectfont} % Times-like (similar to Times New Roman)
\newcommand{\headerfontiii}{\normalfont} % main font already TeX Gyre Pagella via fontspec
\newcommand{\headerfontiv}{\fontfamily{pbk}\selectfont} % Bookman (readable serif)
\newcommand{\headerfontv}{\fontfamily{pag}\selectfont} % Avant Garde-like (similar to Trebuchet MS)
\newcommand{\headerfontvi}{\fontfamily{cmss}\selectfont} % Computer Modern Sans Serif
\newcommand{\headerfontvii}{\fontfamily{qhv}\selectfont} % Quasi-Helvetica (another Arial/Calibri alternative)
\newcommand{\headerfontviii}{\fontfamily{qpl}\selectfont} % Quasi-Palatino (another elegant serif option)
\newcommand{\headerfontix}{\fontfamily{qtm}\selectfont} % Quasi-Times (another Times New Roman alternative)
\newcommand{\headerfontx}{\fontfamily{bch}\selectfont} % Charter (clean serif font)

\begin{document}
\headerfontiii

% Header
\begin{center}
    {\Huge\textbf{王伟钦}}
\end{center}
\vspace{-4.5mm}

\begin{center}
    \small
    {\color{black}{\footnotesize\faPhone*}\hspace*{0.13cm}(+86) 187-6005-3825} |
    \href{mailto:here1swqw@gmail.com}{{\color{black}{\footnotesize\faEnvelope[regular]}\hspace*{0.13cm}here1swqw@gmail.com}} |
    \href{https://mootutu.github.io/}{{\color{black}{\footnotesize\faGlobe}\hspace*{0.13cm}mootutu.github.io}} |
    \href{https://scholar.google.com/citations?user=JonSUasAAAAJ}{{\color{black}{\footnotesize\faGraduationCap}\hspace*{0.13cm}谷歌学术}}
\end{center}
\vspace{-4.5mm}

\begin{center}
    \small
    {\color{black}\footnotesize\faMapMarker*}\hspace*{0.13cm}广东省深圳市南山区深圳大学沧海校区致腾楼 L6-711
\end{center}

% \vspace{-4mm}
\vspace{-2mm}

% \section{\textbf{个人简介}}
% \vspace{1mm}
% \small{
% 本人现为深圳大学 2025 级推免入学的硕士研究生，研究方向为大语言模型后训练、测试时拓展与文本表征。
% }
% \vspace{-1mm}
% \vspace{1mm}

\section{\textbf{教育经历}}
    \vspace{-0.4mm}
    \resumeSubHeadingListStart

    \resumeSubheading
    {深圳大学}{中国，深圳}
    {计算机科学与技术 硕士（推免）}{2025.09 -- 2028.06}
    \resumeItemListStart
      \item 研究方向：大语言模型后训练、测试时拓展与文本表征
      \item 导师：\href{https://ylwangy.github.io/}{王祎乐} 助理教授
    \resumeItemListEnd

    \resumeSubheading
    {福建农林大学}{中国，福州}
    {计算机科学与技术 工学学士}{2021.09 -- 2025.06}

    \resumeSubHeadingListEnd

\vspace{-2mm}
% \vspace{1mm}

\section{\textbf{科研经历}}
\vspace{-0.4mm}
  \resumeSubHeadingListStart
  \resumeSubheading
      {深圳大学 可视计算研究中心（VCC）}{中国，深圳}
      {科研实习生}{2024.10 -- 2025.08}
      \resumeItemListStart
        \item 导师：\href{https://ylwangy.github.io/}{王祎乐} 助理教授
        \item 从事大语言模型推理与自洽性（self-consistency）方法相关研究。
      \resumeItemListEnd
  \resumeSubheading
      {深圳大学}{中国，深圳}
      {《编译原理》课程助教}{2025 春季学期、2026 春季学期}
  \resumePOR{审稿人}{：ICLR (2027)}{}
  \resumeSubHeadingListEnd

\vspace{-2mm}


\section{\textbf{论文发表}}
\vspace{-0.4mm}
\small{(* 同等贡献，\# 通讯作者)}
\vspace{0.5mm}
\resumeItemListStart
  \item Feiyang Li, Shengjing Liu, Qi Zhan, Sijie Cheng, \textbf{Weiqin Wang}, Hongwen Chen, Yuxuan Yang, Wen Wang, Yile Wang, \& Hui Huang. \href{https://arxiv.org/abs/2609.38070}{Probability is Not Enough: Exploring and Counting Divergent Tokens for Reasoning Uncertainty Quantification in LLMs}. \textit{arXiv Preprint}, 2026. [\href{https://github.com/szu-tera/DTC}{code}]
    \begin{itemize}[leftmargin=4.5mm,labelsep=1mm,itemsep=0mm,topsep=0.4mm,parsep=0pt,partopsep=0pt]
      \item \textit{背景：}大模型思维链推理能力增强后，评估与校准其推理置信度、量化答案不确定性愈发重要；现有方法多依赖部分关键 token 的概率，但其内在机制仍不清晰。
      \item \textit{方法：}提出 Divergent Token Confidence（DTC）框架，用同一推理轨迹上两个模型 next-token 分布的 Jensen--Shannon 散度识别“分歧 token”，并统计其数量来估计置信度；无需训练、不改变生成过程，同时支持白盒与黑盒（借助辅助模型）两种设置。
      \item \textit{表现：}在多个模型家族、六个数学基准上优于基于概率与口头表达的基线；白盒下仅凭计数即把平均期望校准误差降至 13.0\%（标准全序列方法为 32.7\%--42.4\%），黑盒下在 DeepSeek-V3.2 上由 32.1\%--40.2\% 降至 13.7\%--16.3\%。
    \end{itemize}
  \item Kaijie Zheng*, \textbf{Weiqin Wang*}, Yile Wang, \& Hui Huang. \href{https://arxiv.org/abs/2605.29751}{DySem: Uncovering Dynamic Semantic Components of Large Language Models for Calculating Semantic Textual Similarity}. \textit{To appear in Proceedings of the 2026 Conference on Empirical Methods in Natural Language Processing (\textbf{EMNLP})}, 2026. [\href{https://github.com/szu-tera/DySem}{code}]
    \begin{itemize}[leftmargin=4.5mm,labelsep=1mm,itemsep=0mm,topsep=0.4mm,parsep=0pt,partopsep=0pt]
      \item \textit{背景：}语义文本相似度（STS）是 NLP 的基础任务。当前基于 LLM 的方法通常固定提取末层隐状态、对所有文本对以固定维度计算相似度，存在两点局限：末层隐状态编码的是通用知识而非纯语义知识；隐层维度普遍很大，为语义表示引入冗余与噪声。
      \item \textit{方法：}提出免训练框架 DySem，基于“文本核心语义在翻译间保持不变”的直觉，通过多语言共识挖掘 LLM 内部更与语义相关的组件，并放弃静态表示空间，为每对文本构建依赖输入的联合语义集，仅在该共享维度子集上计算相似度。
      \item \textit{表现：}在七个标准 STS 基准、十种 LLM 上取得最佳/次佳平均成绩，显著优于 AlignedWVA（配对 t 检验 $p=4.4\times10^{-12}$），指令微调模型上平均提升 4.55 分；平均仅约 1000 维或更少即可取得有竞争力的效果。
    \end{itemize}
  \item Ziyang Chen, Zhenxuan Huang, Yile Wang, \textbf{Weiqin Wang}, Lu Yin, \& Hui Huang. \href{https://aclanthology.org/2026.findings-acl.1858.pdf}{SemPA: Improving Sentence Embeddings of Large Language Models through Semantic Preference Alignment}. \textit{In Findings of the 64th Annual Meeting of the Association for Computational Linguistics (\textbf{ACL Findings})}, 2026. [\href{https://github.com/szu-tera/SemPA}{code}]
    \begin{itemize}[leftmargin=4.5mm,labelsep=1mm,itemsep=0mm,topsep=0.4mm,parsep=0pt,partopsep=0pt]
      \item \textit{背景：}传统句子表征方法基于非生成式预训练模型做 token 级对比学习；近期基于生成式 LLM 的方法要么依赖固定 prompt 模板（未对模型进一步优化、性能受限），要么改动模型结构（改变内部计算机制、损害生成能力）。
      \item \textit{方法：}提出 SemPA，通过语义偏好对齐在保留 LLM 生成能力的同时提升句子表征：在复述生成任务上应用句子级直接偏好优化（DPO），使模型学会判别语义等价的句子，并在 Plackett-Luce 模型框架下建立了 DPO 与对比学习的形式化联系。
      \item \textit{表现：}在语义文本相似度任务及多项 LLM 基准上均取得更好的语义表征，且不牺牲模型固有的生成能力。
    \end{itemize}
  \item \textbf{Weiqin Wang}, Yile Wang, Kehao Chen, \& Hui Huang. \href{https://aclanthology.org/2026.acl-long.1727.pdf}{Beyond Majority Voting: Towards Fine-grained and More Reliable Reward Signal for Test-Time Reinforcement Learning}. \textit{In Proceedings of the 64th Annual Meeting of the Association for Computational Linguistics (\textbf{ACL})}, 2026. [\href{https://github.com/szu-tera/SCOPE}{code}]
    \begin{itemize}[leftmargin=4.5mm,labelsep=1mm,itemsep=0mm,topsep=0.4mm,parsep=0pt,partopsep=0pt]
      \item \textit{背景：}测试时强化学习以多数投票结果作为伪标签，降低对标注数据的依赖，是 RLVR 的互补方向；但该投票策略易引发确认偏差、且奖励稀疏，限制了整体性能。
      \item \textit{方法：}提出框架 SCOPE，将步级置信度融入伪标签估计，优先选择高质量推理路径而非简单计频；并在推理质量与探索多样性之间权衡，把候选输出池动态划分为相互独立的子群，对每个子群重复采样得到局部共识，提供多样化的监督目标以鼓励更广探索。
      \item \textit{表现：}在多种模型与基准上一致优于近期基线；在挑战性的 AIME 2025 上相对提升 13.1\%，在 AMC 上提升 8.1\%。
    \end{itemize}
  \item Zhanyu Shen, Sijie Cheng, Zhicheng Guo, \textbf{Weiqin Wang}, Yile Wang, \& Hui Huang. \href{https://aclanthology.org/2026.findings-acl.1736.pdf}{AnchorMem: Anchored Facts with Associative Contexts for Building Memory in Large Language Models}. \textit{In Findings of the 64th Annual Meeting of the Association for Computational Linguistics (\textbf{ACL Findings})}, 2026. [\href{https://github.com/szu-tera/AnchorMem}{code}]
    \begin{itemize}[leftmargin=4.5mm,labelsep=1mm,itemsep=0mm,topsep=0.4mm,parsep=0pt,partopsep=0pt]
      \item \textit{背景：}LLM 虽在复杂任务上表现优异，但在长期交互中仍需记忆系统来利用历史经验。现有方法（如 A-Mem、Mem0）过度依赖频繁改写来组织交互，而对摘要的过度依赖会稀释关键上下文细节、模糊关键检索特征。
      \item \textit{方法：}受认知科学“普鲁斯特效应”（特定锚点触发整体回忆）启发，提出 AnchorMem，将检索单元与生成上下文解耦：从历史中抽取原子事实作为检索锚点，并保留原始上下文作为不可变上下文；构建带高阶事件连接的关联事件图，把相关事实绑定为共享事件表示。检索时以事实与事件为锚定位相关记忆，再用关联的原始文本块与事件重建上下文。
      \item \textit{表现：}在 LoCoMo 基准上，跨三种闭源与开源模型显著优于基线，在实现细粒度检索的同时保持交互上下文的完整性。
    \end{itemize}
  \item \textbf{Weiqin Wang}, Yile Wang, \& Hui Huang. \href{https://aclanthology.org/2025.findings-acl.744/}{Ranked Voting based Self-Consistency of Large Language Models}. \textit{In Findings of the 63rd Annual Meeting of the Association for Computational Linguistics (\textbf{ACL Findings})}, 2025. [\href{https://github.com/szu-tera/RankedVotingSC}{code}]
    \begin{itemize}[leftmargin=4.5mm,labelsep=1mm,itemsep=0mm,topsep=0.4mm,parsep=0pt,partopsep=0pt]
      \item \textit{背景：}多数投票被视为增强思维链推理的有效方法，即选择不同推理路径中“自洽性”最高的答案；但已有方法每次试验通常只生成单一答案，忽略了其他潜在答案，使这些备选答案在后续投票中被丢弃。
      \item \textit{方法：}提出在每次推理过程中生成排序答案，并在多条响应的排序答案之间进行排序投票，使整体自洽性更可靠；具体采用即时决胜制、Borda 计数与平均倒数秩三种排序投票方法。
      \item \textit{表现：}在六个数据集（三个多选题与三个开放式问答）上，使用开源与闭源 LLM 均优于基线，展现出利用排序答案信息、通过排序投票提升推理性能的潜力。
    \end{itemize}
\resumeItemListEnd

\vspace{-2mm}
\section{\textbf{荣誉奖项}}
    \resumeItemListStart
    \item 国家奖学金（本科生）
    \item 国家奖学金（硕士研究生）
    \item 深圳大学硕士生学业奖学金（特等）
    \item 福建农林大学本科生学业奖学金（一等）
    \item 全国大学生数学建模竞赛二等奖
    \item 美国大学生数学建模竞赛／交叉学科建模竞赛（MCM/ICM）Meritorious Winner（一等奖）
    \resumeItemListEnd

% \vspace{-2mm}
% \section{\textbf{专业技能}}
% \vspace{-0.4mm}
%     \resumeItemListStart
%     \item \textbf{编程与机器学习：}Python, PyTorch, SQL Server, C/C++, LaTeX, Markdown
%     \item \textbf{工具：}Linux, Git, Office
%     \item \textbf{语言：}中文（母语）；英语（CET-4: 532；CET-6: 509；具备学术读写与日常交流能力）
%     \resumeItemListEnd

\end{document}
